\documentclass[letterpaper]{article}

\usepackage[letterpaper,margin=1in]{geometry}

% Compile from the repository root so this path resolves to the pinned WSP
% submodule.
\input{wsp/testing/test-case-library.tex}

\begin{document}
\def\TCID{TC-0005}
\def\TCTitle{PDS programs in WFC}
\def\TCPurpose{Verify that WFC runs PDS 7.1 programs on the runtime, in the
interpreter and as native code: the statements and functions of the text
screen, PRINT on the screen, INKEY\$ from the standard input, error 5, and
VB6 modules that use the same words as names.}
\def\TCPriority{\TCPriorityHigh}

\def\TCRequirementRef{REQ-0004, REQ-0005, REQ-0006, REQ-0007, REQ-0008,
REQ-0009, REQ-0010, REQ-0011}
\def\TCDesignRef{WFC \texttt{architecture/adr-0011-pds-text-programs.md}.}
\def\TCFeatureRef{\texttt{tests/corpus/*.bas}; in WFC
\texttt{src/interpreter/pds\_screen.cpp},
\texttt{src/interpreter/pds\_statements.cpp}.}
\def\TCTestTechnique{Scenario tests: each program prints, after its plain
output, the text of its final screen, which is compared with the expected
text.}

\def\TCPreconditions{A \texttt{wfc.exe} built with this library; CMake is
given \texttt{-DWFC\_EXECUTABLE=<wfc.exe>}. Without it the tests are skipped.
Inside WFC the programs run as \texttt{TC-MP0002-corpus-NNN-name} and
\texttt{TC-MP0002-corpus-NNN-name-native}.}
\def\TCEnvironment{Windows. The standard output is redirected, so the screen
is kept in memory and its text printed at the end. CTest runs each program
with a 120 second timeout.}
\def\TCAssumptions{The expected texts follow REQ-0001 and REQ-0004; they were
not taken from BRUN71 itself.}

\def\TCInputData{The programs \texttt{300-pds-screen},
\texttt{301-pds-inkey} (with \texttt{301-pds-inkey.keys} as standard input),
\texttt{302-pds-errors}, \texttt{303-pds-names}, \texttt{304-pds-scroll},
\texttt{305-pds-input} (with \texttt{305-pds-input.keys}),
\texttt{306-pds-using}, \texttt{307-pds-sound-keys} (with
\texttt{307-pds-sound-keys.keys}), \texttt{308-pds-sleep-api},
\texttt{309-pds-on-key} (with \texttt{309-pds-on-key.keys}),
\texttt{310-pds-write-screen}, \texttt{311-pds-error-exit} and the form
projects \texttt{320-text-form.vbp}, \texttt{321-text-input.vbp},
\texttt{322-text-lists.vbp} and \texttt{323-text-combo-open.vbp} (with
their \texttt{.keys} files).}
\def\TCInitialState{None.}

\def\TCProcedure{\begin{enumerate}
  \item Run each program with \texttt{wfc --interpret} (CTest
    \texttt{TC-0005-NNN-name}) and as native code
    (\texttt{TC-0005-NNN-name-native}).
  \item Compare the standard output with \texttt{NNN-name.expected}, ignoring
    the line ending and the blanks around the whole text.
\end{enumerate}}
\def\TCExpectedResult{Every program exits with status 0 (or the status of
its \texttt{.exitcode} file, 1 for \texttt{311-pds-error-exit}) and prints
exactly its expected text in both modes.}
\def\TCPostConditions{\texttt{303-pds-names} and \texttt{306-pds-using}
delete the files they write.}

\TCTemplate

\end{document}
